% 第1回　ガイダンス（記入例・模範解答）
\session{1}{1}{9月30日（水）}{ガイダンス}{}{}

\goals{
  \item 授業の二部構成と評価の仕組みを説明できる
  \item この授業でのAI利用のルールを説明できる
  \item ハルシネーションがなぜ起きるかを理解し、出力を確かめる方法を挙げられる
}

\work{確認}{授業の全体像と評価}
\begin{wstab}{colspec={Q[l,wd=40mm,m] X[l,m]},row{1-Z}={ht=9mm}}
  \hc{前半（第1〜5回）} & \af{探究の基礎。AIを使って問いを立て、先行研究を探して読み、リサーチ・クエスチョンを確定する} \\
  \hc{後半（第6〜13回）} & \af{実践プロジェクト。企業の課題（カーディーラーの店舗設計）に、グループで提案をつくる} \\
  \hc{成績評価の内訳} & 提出物 \underline{\makebox[12mm]{\ans{50}}}\,\%　／　発表 \underline{\makebox[12mm]{\ans{50}}}\,\% \\
  \hc{主な提出物} & \af{前半：リサーチ・クエスチョン確定レポート／後半：店舗設計提案／最後：振り返りレポート（いずれもAI活用記録を添付）} \\
\end{wstab}

\work{確認}{この授業でのAI利用のルール（自分の言葉で3つ）}
\alines{3}{① 大学が提供する Copilot を使い、企業の非公開情報や個人情報は入力しない。\par
② どこでAIを使ったかを記録し、提出物に「AI活用記録」として明記する。\par
③ 出力は素材として扱い、事実・数値・文献は一次資料で確かめてから使う。}

\work[12mm]{準備}{Copilot の動作確認}
\begin{achecks}
  \item 大学のアカウントで Copilot にサインインできた
  \item 簡単な質問をして、回答が返ってくることを確かめた
\end{achecks}

\work[70mm]{ミニワーク}{AIに嘘をつかせる}
\begin{promptbox}[ハルシネーションを体験する]
名古屋市立大学経済学部の歴史について、参考文献を3つ挙げて300字で説明してください。
\end{promptbox}
\instr{(1) AIが挙げた参考文献が実在するか、図書館の蔵書検索や CiNii Research で確かめる。}
\begin{wstabh}{colspec={X[3,l,m] Q[c,wd=14mm,m] X[2,l,m]},row{2-Z}={ht=10mm}}
  AIが挙げた文献（著者・題名・年） & 実在 ○× & 確かめた場所・わかったこと \\
  \af{『名古屋市立大学経済学部五十年史』（編者名・発行年つきで提示）} & \ans{×} & \af{蔵書検索・CiNii とも該当なし。題名も編者も確認できない} \\
  \af{「公立大学における経済学教育の展開」（学会誌の論文として提示）} & \ans{×} & \af{CiNii・J-STAGE で該当なし。掲載誌は実在するが、その号にない} \\
  \af{大学公式サイトの「沿革」ページ} & \ans{○} & \af{実在。ただしAIの説明にある設立年と食い違う} \\
\end{wstabh}

\instr{(2) 自分が詳しく知っている話題（出身地、部活、趣味など）について聞き、内容の正誤を判定する。}
\begin{wstabh}{colspec={X[2,l,m] X[3,l,m] Q[c,wd=14mm,m] X[2,l,m]},row{2-Z}={ht=12mm}}
  尋ねたこと & AIの答え（要点） & 正誤 ○△× & どこが違ったか \\
  \af{出身地の夏祭りの由来} & \af{起源の時代と主催団体を断定的に説明} & \ans{△} & \af{主催団体は正しいが、起源の時代は地元資料と違う} \\
  \af{高校の部活（吹奏楽）の大会の仕組み} & \af{地区大会から全国大会までの流れ} & \ans{○} & \af{おおむね正しい。細かな日程は年によって違う} \\
  \af{地元の駅の開業年} & \af{具体的な年を1つ挙げた} & \ans{×} & \af{鉄道会社の公式サイトの年と10年以上ずれていた} \\
\end{wstabh}

\work{まとめ}{気づいたこと・「出力を使う前のチェック」}
\instr{共有で出た意見も含め、AIの出力を使う前に確かめるべきことを自分の言葉で整理する。}
\alines{5}{AIは、知らないことも自信のある口調で答える。文献は「それらしい題名」で作られることがある。\par
・固有名詞（書名・人名・組織名）は、蔵書検索や CiNii で実在を確かめる\par
・数値や年は、出典にあたって同じ値かを確かめる\par
・よく知っている話題で一度試し、どの程度間違えるかの感覚を持っておく\par
・自分で説明できない内容は使わない}

\begin{nexttime}
  \item 関心のあるテーマを3つ書き出す（第2回のブレストで使う）。「なぜ気になるのか」を一言添える。
\end{nexttime}
\begin{wstabh}{colspec={Q[c,wd=8mm,m] X[2,l,m] X[3,l,m]},row{2-Z}={ht=11mm}}
  & 関心のあるテーマ & なぜ気になるのか \\
  1 & \af{地方の自動車販売と若い世代} & \af{実家の近くのディーラーがいつも空いていて、同世代が車を買う話も聞かないから} \\
  2 & \af{大学周辺の飲食店のキャッシュレス化} & \af{現金のみの店とキャッシュレスの店で、客層が違うように見えるから} \\
  3 & \af{「推し活」にお金を使う理由} & \af{周りの友人の支出の中で、推し活の割合が大きいと感じるから} \\
\end{wstabh}

\begin{point}
  \item 評価の内訳（提出物50\%・発表50\%）と、提出物には必ずAI活用記録を付けることを、この回で確実に押さえさせる。
  \item AI利用のルールは「大学のCopilotを使う」「記録する」「検証する」の3点が書けていればよい。入力してはいけない情報（個人情報・企業の非公開情報）に触れていれば加点的に評価する。
  \item ミニワークでは、存在しない文献が「実在しそうな体裁」で出ることを体験させるのが目的。AIの出力は回ごとに変わるので、表の記入例どおりになるとは限らない。実在しない文献が1件も出なかった学生には、その結果自体を記録させ、「今回は出なかった」ことも検証の成果だと伝える。
  \item 関心テーマは、理由まで書けているかを見る。「なぜ気になるか」が第2回で問いに育てる手がかりになる。
\end{point}
